% =====================================================================
%  ANEXO A. Especificación de casos de prueba
% =====================================================================
\section{Especificación de casos de prueba}
\label{anx:casos}

\begin{notanormativa}[Qué es y qué no es un caso de prueba]
Un caso de prueba describe \textbf{qué debería ejecutarse y qué debería ocurrir}. Es un
artefacto de diseño y se mantiene estable entre ejecuciones. Lo que \emph{sucedió} al ejecutarlo
—resultado real, veredicto, evidencia— no se escribe aquí: se registra en el
\ref{anx:ejecuciones}. Esa separación es la que permite ejecutar el mismo caso muchas veces sin
sobrescribir su historial.
\end{notanormativa}

\begin{instruccion}
Reproduzca la ficha siguiente una vez por cada caso de prueba. Un caso queda derivado de una
condición de prueba (apartado~\ref{sec:diseno}, Cuadro~\ref{tab:condiciones}), que a su vez
procede de la base de prueba y
de un riesgo: por eso la ficha pide esas tres claves. El Cuadro~\ref{tab:indice-casos} sirve
como índice de todos los casos especificados.

Si mantiene los casos en una herramienta externa, conserve al menos los mismos campos e indique
su ubicación en el Cuadro~\ref{tab:configuracion}.
\end{instruccion}

\begin{observacion}[Anexo sin formato]
\obsproblema{El anexo equivalente de la plantilla anterior decía únicamente «incluir los casos
de prueba aplicados», sin campos ni estructura, y el apartado que lo invocaba mezclaba el diseño
del caso con el resultado de su ejecución.}
\obscambio{Ficha completa con base, riesgo, condición, técnica, elemento de cobertura,
precondiciones, datos, pasos, resultado esperado, postcondiciones y limpieza, más un índice de
casos. Los campos de resultado se trasladan al \ref{anx:ejecuciones}.}
\obsfuente{OBS §2 y §6}
\end{observacion}

\subsection*{Ficha de caso de prueba}

\begin{xltabular}{\textwidth}{|E{4.4cm}|Y|}
\caption{Formato de especificación de un caso de prueba.}\label{tab:ficha-caso}\\ \hline
\endfirsthead
\multicolumn{2}{@{}l@{}}{\footnotesize\itshape\tablename~\thetable{} (continuación)}\\ \hline
\endhead
\multicolumn{2}{@{}r@{}}{\footnotesize\itshape continúa en la página siguiente}\\
\endfoot
\endlastfoot
  Identificador del caso        & \campo{CP-001} \\ \hline
  Versión del caso              & \campo{v1} \\ \hline
  Objeto de prueba              & \campo{OP-01 (Cuadro~\ref{tab:objetos})} \\ \hline
  Base de prueba                & \campo{BP-01 (Cuadro~\ref{tab:base-prueba})} \\ \hline
  Riesgo asociado               & \campo{RP-01 (Cuadro~\ref{tab:riesgos-producto})} \\ \hline
  Condición de prueba           & \campo{COND-01 (Cuadro~\ref{tab:condiciones})} \\ \hline
  Objetivo del caso             & \campo{Qué pretende evidenciar este caso en particular} \\ \hline
  Nivel de prueba               & \campo{Componente / integración / sistema / aceptación} \\ \hline
  Técnica aplicada              & \campo{Cuadro~\ref{tab:tecnicas}} \\ \hline
  Elemento(s) de cobertura      & \campo{Qué partición, valor límite, decisión o transición ejercita} \\ \hline
  Prioridad                     & \campo{Alta / media / baja, según el riesgo} \\ \hline
  Precondiciones                & \campo{Estado que debe tener el sistema antes de ejecutar} \\ \hline
  Datos de entrada              & \campo{Valores concretos y conjunto de datos (DAT-xx)} \\ \hline
  Acciones / pasos              & \campo{1. …\quad 2. …\quad 3. …} \\ \hline
  Resultado esperado            & \campo{Qué debe ocurrir; indique el oráculo del que se deriva} \\ \hline
  Postcondiciones               & \campo{Estado en que queda el sistema tras la ejecución} \\ \hline
  Limpieza requerida            & \campo{Qué debe deshacerse o restaurarse después} \\ \hline
  Dependencias                  & \campo{Otros casos que deban ejecutarse antes} \\ \hline
  Autor y fecha                 & \campo{Nombre — dd/mm/aaaa} \\ \hline
\end{xltabular}

El Cuadro~\ref{tab:ficha-caso} es el formato que debe replicarse por cada caso. Sus campos de
base, riesgo y condición son los que alimentan la matriz de trazabilidad del
\ref{anx:trazabilidad}; el campo «resultado esperado» es el que se referencia —no se copia—
desde cada ejecución.

\subsection*{Índice de casos de prueba especificados}

\begin{longtable}{|C{1.6cm}|L{4.0cm}|C{1.6cm}|C{1.6cm}|C{1.9cm}|C{1.6cm}|}
\caption{Índice de casos de prueba especificados.}
\label{tab:indice-casos} \\
\hline
\filaenc \textbf{ID} & \textbf{Título del caso} & \textbf{Objeto} & \textbf{Condición} & \textbf{Técnica} & \textbf{Prioridad} \\ \hline
\endfirsthead
\hline
\filaenc \textbf{ID} & \textbf{Título del caso} & \textbf{Objeto} & \textbf{Condición} & \textbf{Técnica} & \textbf{Prioridad} \\ \hline
\endhead
\hline \multicolumn{6}{|r|}{\footnotesize\textit{continúa en la página siguiente}} \\ \hline
\endfoot
\hline
\endlastfoot
\campo{CP-001} & \campo{Título} & \campo{OP-01} & \campo{COND-01} & \campo{Valores límite} & \campo{Alta} \\ \hline
 & & & & & \\ \hline
 & & & & & \\ \hline
 & & & & & \\ \hline
\end{longtable}

El Cuadro~\ref{tab:indice-casos} permite ver de un vistazo el conjunto diseñado y detectar
condiciones sin casos o técnicas declaradas que no se usaron.
